\documentclass[a4paper]{article}
\usepackage[fontset=fandol]{ctex}
\usepackage{amsmath, amssymb}
\usepackage{graphicx}
\usepackage[margin=2.5cm]{geometry}
\usepackage[most]{tcolorbox}
\usepackage{etoolbox}
\usepackage{listings}
\usepackage{hyperref}
\usepackage{booktabs}
\usepackage{subcaption}
\usepackage{float}
\usepackage{tikz}
\IfFileExists{pgfplots.sty}{\usepackage{pgfplots}\pgfplotsset{compat=1.18}}{}
\newtcolorbox{knowledgebox}[1]{enhanced, colback=blue!5!white, colframe=blue!75!black, colbacktitle=blue!75!black, coltitle=white, fonttitle=\bfseries, title=#1, attach boxed title to top left={yshift=-2mm, xshift=2mm}, boxrule=1pt, sharp corners}
\newtcolorbox{importantbox}[1]{enhanced, colback=yellow!10!white, colframe=yellow!80!black, colbacktitle=yellow!80!black, coltitle=black, fonttitle=\bfseries, title=#1, sharp corners}
\newtcolorbox{warningbox}[1]{enhanced, colback=red!5!white, colframe=red!75!black, colbacktitle=red!75!black, coltitle=white, fonttitle=\bfseries, title=#1, sharp corners}
\lstset{language=Python,basicstyle=\ttfamily\small,keywordstyle=\color{blue},stringstyle=\color{red!60!black},commentstyle=\color{green!60!black},breaklines=true,frame=single,numbers=left,numberstyle=\tiny\color{gray},captionpos=b,extendedchars=false}
\newcommand{\notetitle}{Introduction of Speech Recognition}
\newcommand{\noteauthors}{Codex 基于课程页、公开视频字幕、coverage ledger 与补充官方材料整理}
\newcommand{\notedate}{\today}
\newcommand{\videochannel}{WAVLab}
\newcommand{\videopublishdate}{2023-08-30}
\newcommand{\videoduration}{unknown}
\newcommand{\videourl}{https://www.youtube.com/watch?v=thEKvvFwMvQ}
\newcommand{\videocoverpath}{cover.jpg}
\begin{document}
\begin{titlepage}
\centering
{\Large 课程笔记\par}
\vspace{1.2cm}
{\huge\bfseries \notetitle\par}
\vspace{0.8cm}
{\large \noteauthors\par}
\vspace{0.3cm}
{\large \notedate\par}
\vspace{1.2cm}
\ifdefempty{\videocoverpath}{{\small 公开视频封面不可用。\par}}{\includegraphics[width=0.82\textwidth,height=0.45\textheight,keepaspectratio]{\videocoverpath}\par}
\vfill
\begin{tcolorbox}[width=0.9\textwidth, colback=black!2!white, colframe=black!60, sharp corners]
\textbf{视频作者/频道}：\videochannel\par
\textbf{发布日期}：\videopublishdate\par
\textbf{视频时长}：\videoduration\par
\textbf{视频链接}：\href{\videourl}{\nolinkurl{\videourl}}
\end{tcolorbox}
\end{titlepage}
\tableofcontents
\newpage
\section{2.0 课程主题总览}
根据课程页主题、公开视频字幕以及本讲的 coverage ledger，2.0 课程主题总览 这一部分主要围绕 Introduction of Speech Recognition | Evaluation metric | How to transcribe speech | Databases 展开。
写作上保持 coverage-first 原则：先把老师在这一阶段真正讲过的问题设定、模型部件、训练或推理逻辑讲清，再压缩成适合复习的结构。

\begin{figure}[H]
\centering
\includegraphics[width=0.92\textwidth]{figures/coverage_map.png}
\caption{本讲 coverage map：按 coverage\_units.jsonl 归纳出的核心主题与章节映射。}
\end{figure}

\begin{figure}[H]
\centering
\includegraphics[width=0.92\textwidth]{figures/timeline.png}
\caption{本讲 timeline：根据字幕时间范围归纳的主题推进顺序。}
\end{figure}

\subsection{Introduction of Speech Recognition}
从证据位置看，这一部分对应的课程证据区间是 2023-08-30。课程在这里强调的核心内容可以概括为：Introduction of Speech Recognition | Evaluation metric | How to transcribe speech | Databases。
这一部分把语音识别系统的好坏从“听起来差不多”转成可重复计算的误差统计，这也是后续比较不同模型与搜索策略时的统一标尺。

记录系统中的附加说明指出：Course-page schedule entry is restated in the chapter introduction and section roadmap.

在本章结构中，这一单元被并入“2.0 课程主题总览”，目的是把同一阶段的定义、机制和权衡集中讲清。

\subsection{本章小结}
2.0 课程主题总览 这一部分把 Introduction of Speech Recognition 中相邻的教学单元串成了一条连续叙事：既交代了这一阶段要解决的问题，也交代了方法、推理或工程权衡为什么要这样组织。

\section{2.1 评价指标（evaluation metric）}
根据课程页主题、公开视频字幕以及本讲的 coverage ledger，2.1 评价指标（evaluation metric） 这一部分主要围绕 Evaluation metric | Word error rate | Edit distance intuition 展开。
写作上保持 coverage-first 原则：先把老师在这一阶段真正讲过的问题设定、模型部件、训练或推理逻辑讲清，再压缩成适合复习的结构。

\subsection{Evaluation metric}
从证据位置看，这一部分对应的课程证据区间是 00:00:00,080 -- 00:07:19,919。课程在这里强调的核心内容可以概括为：Evaluation metric | Word error rate | Edit distance intuition。
这一部分把语音识别系统的好坏从“听起来差不多”转成可重复计算的误差统计，这也是后续比较不同模型与搜索策略时的统一标尺。

记录系统中的附加说明指出：WER and edit distance are expanded with formula and symbol explanation.

在本章结构中，这一单元被并入“2.1 评价指标（evaluation metric）”，目的是把同一阶段的定义、机制和权衡集中讲清。
\[
\mathrm{WER} = \frac{S + D + I}{N}
\]

符号说明如下：
\begin{itemize}
\item $S$：substitutions，替换错误数；
\item $D$：deletions，删除错误数；
\item $I$：insertions，插入错误数；
\item $N$：参考转写中的词数。
\end{itemize}

\subsection{本章小结}
2.1 评价指标（evaluation metric） 这一部分把 Introduction of Speech Recognition 中相邻的教学单元串成了一条连续叙事：既交代了这一阶段要解决的问题，也交代了方法、推理或工程权衡为什么要这样组织。

\section{2.2 转写规范（transcription policy）}
根据课程页主题、公开视频字幕以及本讲的 coverage ledger，2.2 转写规范（transcription policy） 这一部分主要围绕 How to transcribe speech | Orthographic normalization | Boundary decisions 展开。
写作上保持 coverage-first 原则：先把老师在这一阶段真正讲过的问题设定、模型部件、训练或推理逻辑讲清，再压缩成适合复习的结构。

\subsection{How to transcribe speech}
从证据位置看，这一部分对应的课程证据区间是 00:07:16,120 -- 00:13:45,399。课程在这里强调的核心内容可以概括为：How to transcribe speech | Orthographic normalization | Boundary decisions。
课程强调，转写规范并不是纯粹的标注细节，而是直接决定训练目标、评测方式与跨语料可比性的建模前提。

记录系统中的附加说明指出：Orthographic normalization, punctuation, and boundary conventions are covered as transcription policy.

在本章结构中，这一单元被并入“2.2 转写规范（transcription policy）”，目的是把同一阶段的定义、机制和权衡集中讲清。

\subsection{本章小结}
2.2 转写规范（transcription policy） 这一部分把 Introduction of Speech Recognition 中相邻的教学单元串成了一条连续叙事：既交代了这一阶段要解决的问题，也交代了方法、推理或工程权衡为什么要这样组织。

\section{2.3 语音数据库（databases）}
根据课程页主题、公开视频字幕以及本讲的 coverage ledger，2.3 语音数据库（databases） 这一部分主要围绕 Databases | Corpus design | Train/dev/test splits 展开。
写作上保持 coverage-first 原则：先把老师在这一阶段真正讲过的问题设定、模型部件、训练或推理逻辑讲清，再压缩成适合复习的结构。

\subsection{Databases}
从证据位置看，这一部分对应的课程证据区间是 00:13:42,680 -- 00:21:03,559。课程在这里强调的核心内容可以概括为：Databases | Corpus design | Train/dev/test splits。
从教材视角看，语料库设计决定了系统能学到什么分布，也决定了实验结果是否能被正确解释。

记录系统中的附加说明指出：Corpus design, split hygiene, and scale tradeoffs are covered.

在本章结构中，这一单元被并入“2.3 语音数据库（databases）”，目的是把同一阶段的定义、机制和权衡集中讲清。

\subsection{本章小结}
2.3 语音数据库（databases） 这一部分把 Introduction of Speech Recognition 中相邻的教学单元串成了一条连续叙事：既交代了这一阶段要解决的问题，也交代了方法、推理或工程权衡为什么要这样组织。

\section{2.4 语料比较与标注一致性}
根据课程页主题、公开视频字幕以及本讲的 coverage ledger，2.4 语料比较与标注一致性 这一部分主要围绕 Corpus comparison | Annotation consistency | Subjective judgments 展开。
写作上保持 coverage-first 原则：先把老师在这一阶段真正讲过的问题设定、模型部件、训练或推理逻辑讲清，再压缩成适合复习的结构。

\subsection{Corpus comparison}
从证据位置看，这一部分对应的课程证据区间是 00:20:59,320 -- 00:27:26,559。课程在这里强调的核心内容可以概括为：Corpus comparison | Annotation consistency | Subjective judgments。
从教材视角看，语料库设计决定了系统能学到什么分布，也决定了实验结果是否能被正确解释。

记录系统中的附加说明指出：Comparison criteria are written up as annotation consistency and subjective judgment caveats.

在本章结构中，这一单元被并入“2.4 语料比较与标注一致性”，目的是把同一阶段的定义、机制和权衡集中讲清。

\subsection{本章小结}
2.4 语料比较与标注一致性 这一部分把 Introduction of Speech Recognition 中相邻的教学单元串成了一条连续叙事：既交代了这一阶段要解决的问题，也交代了方法、推理或工程权衡为什么要这样组织。

\section{2.5 录制与数据收集}
根据课程页主题、公开视频字幕以及本讲的 coverage ledger，2.5 录制与数据收集 这一部分主要围绕 Speech data collection | Spontaneous speech | Recording conditions 展开。
写作上保持 coverage-first 原则：先把老师在这一阶段真正讲过的问题设定、模型部件、训练或推理逻辑讲清，再压缩成适合复习的结构。

\subsection{Speech data collection}
从证据位置看，这一部分对应的课程证据区间是 00:27:22,880 -- 00:33:40,279。课程在这里强调的核心内容可以概括为：Speech data collection | Spontaneous speech | Recording conditions。
从课程组织上看，这一部分承担的是把局部概念嵌回整条识别流水线中的作用，避免读者把公式或模块孤立地记忆。

记录系统中的附加说明指出：Spontaneous speech and recording conditions are described in the data-collection section.

在本章结构中，这一单元被并入“2.5 录制与数据收集”，目的是把同一阶段的定义、机制和权衡集中讲清。

\subsection{本章小结}
2.5 录制与数据收集 这一部分把 Introduction of Speech Recognition 中相邻的教学单元串成了一条连续叙事：既交代了这一阶段要解决的问题，也交代了方法、推理或工程权衡为什么要这样组织。

\section{2.6 基准、工具链与实践注意事项}
根据课程页主题、公开视频字幕以及本讲的 coverage ledger，2.6 基准、工具链与实践注意事项 这一部分主要围绕 Resource scale | Toolkits | Benchmarking and practical caveats 展开。
写作上保持 coverage-first 原则：先把老师在这一阶段真正讲过的问题设定、模型部件、训练或推理逻辑讲清，再压缩成适合复习的结构。

\subsection{Resource scale}
从证据位置看，这一部分对应的课程证据区间是 00:33:36,320 -- 00:40:31,000。课程在这里强调的核心内容可以概括为：Resource scale | Toolkits | Benchmarking and practical caveats。
从课程组织上看，这一部分承担的是把局部概念嵌回整条识别流水线中的作用，避免读者把公式或模块孤立地记忆。

记录系统中的附加说明指出：Toolkits, benchmarking, and practical cautions are covered, with CS224S mentioned only as an extension where relevant.

在本章结构中，这一单元被并入“2.6 基准、工具链与实践注意事项”，目的是把同一阶段的定义、机制和权衡集中讲清。
\begin{knowledgebox}{CS224S 补充}
Toolkits, benchmarking, and practical cautions are covered, with CS224S mentioned only as an extension where relevant.
\end{knowledgebox}

\subsection{本章小结}
2.6 基准、工具链与实践注意事项 这一部分把 Introduction of Speech Recognition 中相邻的教学单元串成了一条连续叙事：既交代了这一阶段要解决的问题，也交代了方法、推理或工程权衡为什么要这样组织。

\section{总结与延伸}
Introduction of Speech Recognition 这一讲在整门《Speech Recognition and Understanding》课程中的位置，是把 Introduction of Speech Recognition 讲成一个可复用的建模模块。
从教材化角度看，本讲最重要的不是记住单一术语，而是理解它如何嵌入完整的 automatic speech recognition pipeline：输入语音如何被表示、模型如何被训练、解码如何得到最终转写，以及这些设计分别带来什么假设和权衡。
由于公开源以课程页和公开视频字幕为主、缺少官方 slide PDF，本章在正文里尽量把术语、方法和工程语境解释完整，并把缺失项留在 omission 和 manifest 里，而不是凭空补写。
\end{document}
